\documentclass[10pt,
% a4paper,
%twocolumn,
fleqn,
%landscape, 
% papersize,
dvipdfmx,
uplatex
]{jsarticle}



\def\maru#1{\textcircled{\scriptsize#1}}%丸囲み番号

% \RequirePackage[2020/09/30]{platexrelease}

%太字設定
\usepackage[deluxe]{otf}

\usepackage{emathEy}

\usepackage[g]{esvect}

%大きな文字
\usepackage{fix-cm}

%定理環境
\usepackage{emathThm}
%\theoremstyle{boxed}
\theorembodyfont{\normalfont}
\newtheorem{Question}{問題}[subsection]
\newtheorem{Q}{}[subsection]
\newtheorem{question}[Question]{}
\newtheorem{quuestion}{}[subsection]

%セクション，大問番号のデザイン
\renewcommand{\labelenumi}{(\arabic{enumi})}
\renewcommand{\theenumii}{\alph{enumii})}
\renewcommand{\thesection}{第\arabic{section}章}

%用紙サイズの詳細設定
\usepackage{bxpapersize}
\papersizesetup{size={80mm,45mm}}
\usepackage[top=0.7zw,bottom=0truemm,left=3truemm,right=133truemm]{geometry}
\usepackage[dvipdfmx]{graphicx}

%余白など
\usepackage{setspace} % 行間
\setlength{\mathindent}{1zw}
\setlength\parindent{0pt}


%色カラーに関する設定
\usepackage{color}
\definecolor{shiro}{rgb}{0.95703125,0.87109375,0.7421875}
\definecolor{kin}{rgb}{0.95703125,0.87109375,0.7421875}
\definecolor{orange}{rgb}{1,0.7,0.2}
\definecolor{bradorange}{rgb}{1,0.5,0}
\definecolor{pink}{rgb}{0.9176,0.5686,0.5960}
\definecolor{mizu}{rgb}{0.6156,0.8,0.9955}
\color{kin}
% \pagecolor{hukamido}

\usepackage{at}%図の配置
% \usepackage{wallpaper}

\begin{document}



\at(0cm,0cm){\includegraphics[width=8cm,bb=0 0 1920 1080]{./youtube/thumbnails/templates/smart_background/三角関数.jpeg}}
{\color{orange}\bf\boldmath\Large\underline{$2$倍角，$3$倍角の公式の導出}}\vspace{0.3zw}

\bf\boldmath 

\large
\hspace{0.5zw}$\sin 2\alpha =2\sin \alpha \cos \alpha$\\
\hspace{0.5zw}$\cos 2\alpha =1-2\sin ^2\alpha =2\cos ^2-1$\\
\hspace{0.5zw}$\tan 2\alpha =\bunsuu{{2\tan \alpha }}{{1-\tan ^2\alpha }}$\vspace{0.4zw}\\
\hspace{0.5zw}$\sin 3\alpha =3\sin \alpha -4\sin ^3\alpha$\\
\hspace{0.5zw}$\cos 3\alpha =-3\cos \alpha +4\sin ^3\alpha$\\

\at(7.0cm,0.2cm){\small\color{bradorange}$\overset{\text{三角関数}}{\text{基本}}$}


\newpage



\at(0cm,0cm){\includegraphics[width=8cm,bb=0 0 1920 1080]{./youtube/thumbnails/templates/smart_background/三角関数.jpeg}}
{\color{orange}\bf\boldmath\LARGE\underline{$\tan$の加法定理の導出}}\vspace{0.3zw}

\normalsize
\bf\boldmath

\Large
\vspace{0.5zw}
\hspace{0.2zw}$\tan \left(\alpha +\beta \right)=\bunsuu{{\tan \alpha +\tan \beta }}{{1-\tan \alpha \tan \beta }}\vspace{0.3zw}$\vspace{0.5zw}\\
\hspace{0.2zw}$\tan \left(\alpha -\beta \right)=\bunsuu{{\tan \alpha -\tan \beta }}{{1+\tan \alpha \tan \beta }}\vspace{0.3zw}$


\at(7.0cm,0.2cm){\small\color{bradorange}$\overset{\text{三角関数}}{\text{基本}}$}


\newpage



\at(0cm,0cm){\includegraphics[width=8cm,bb=0 0 1920 1080]{./youtube/thumbnails/templates/smart_background/三角関数.jpeg}}
{\color{orange}\bf\boldmath\LARGE\underline{$\sin$,$\cos$の加法定理の導出}}\vspace{0.5zw}

\large
\bf\boldmath 
 $\sin \left(\alpha +\beta \right)=\sin \alpha \cos \beta +\cos \alpha \sin \beta$\vspace{0.3zw}\\
 $\sin \left(\alpha -\beta \right)=\sin \alpha \cos \beta -\cos \alpha \sin \beta$\vspace{0.3zw}\\
 $\cos \left(\alpha +\beta \right)=\cos \alpha \cos \beta -\sin \alpha \sin \beta$\vspace{0.3zw}\\
 $\cos \left(\alpha -\beta \right)=\cos \alpha \cos \beta +\sin \alpha \sin \beta$\vspace{0.3zw}\\

\at(7.0cm,0.2cm){\small\color{bradorange}$\overset{\text{三角関数}}{\text{基本}}$}


\newpage



\at(0cm,0cm){\includegraphics[width=8cm,bb=0 0 1920 1080]{./youtube/thumbnails/templates/smart_background/三角関数.jpeg}}
{\color{orange}\bf\boldmath\Large\underline{三角関数の性質 全パターン導出}}\vspace{0.2zw}

\scriptsize
\bf\boldmath

$\begin{array}{ll}

\left\{\begin{array}{l}
\hspace{-0.5zw}\sin \left(\theta +2n\pi \right)=\sin \theta \vspace{0.3zw}\\
\hspace{-0.5zw}\cos \left(\theta +2n\pi \right)=\cos \theta \vspace{0.3zw}\\
\hspace{-0.5zw}\tan \left(\theta +2n\pi \right)=\tan \theta \vspace{0.3zw}
\end{array}\right.&

\left\{\begin{array}{l}
\hspace{-0.5zw}\sin \left(-\theta \right)=-\sin \theta \vspace{0.3zw}\\
\hspace{-0.5zw}\cos \left(-\theta \right)=\cos \theta \vspace{0.3zw}\\
\hspace{-0.5zw}\tan \left(-\theta \right)=-\tan \theta \vspace{0.3zw}
\end{array}\right.\vspace{0.5zw}\\

\left\{\begin{array}{l}
\hspace{-0.5zw}\sin \left(\theta +\frac{\pi }{2}\right)=\cos \theta \vspace{0.3zw}\\
\hspace{-0.5zw}\cos \left(\theta +\frac{\pi }{2}\right)=-\sin \theta \vspace{0.3zw}\\
\hspace{-0.5zw}\tan \left(\theta +\frac{\pi }{2}\right)=-\frac{1}{\tan\theta} \vspace{0.3zw}
\end{array}\right.&

\left\{\begin{array}{l}
\hspace{-0.5zw}\sin \left(\frac{\pi }{2}-\theta\right)=\cos \theta \vspace{0.3zw}\\
\hspace{-0.5zw}\cos \left(\frac{\pi }{2}-\theta\right)=\sin \theta \vspace{0.3zw}\\
\hspace{-0.5zw}\tan \left(\frac{\pi }{2}-\theta\right)=\frac{1}{\tan\theta} \vspace{0.3zw}
\end{array}\right.\vspace{0.5zw}\\

\left\{\begin{array}{l}
\hspace{-0.5zw}\sin \left(\theta +\pi \right)=-\sin \theta \vspace{0.3zw}\\
\hspace{-0.5zw}\cos \left(\theta +\pi \right)=-\cos \theta \vspace{0.3zw}\\
\hspace{-0.5zw}\tan \left(\theta +\pi \right)=\tan \theta \vspace{0.3zw}
\end{array}\right.&

\left\{\begin{array}{l}
\hspace{-0.5zw}\sin \left(\pi-\theta \right)=\sin \theta \vspace{0.3zw}\\
\hspace{-0.5zw}\cos \left(\pi-\theta\right)=-\cos \theta \vspace{0.3zw}\\
\hspace{-0.5zw}\tan \left(\pi-\theta \right)=-\tan \theta \vspace{0.3zw}
\end{array}\right.\vspace{0.5zw}\\




\end{array}$

\at(7.0cm,0.2cm){\small\color{bradorange}$\overset{\text{三角関数}}{\text{基本}}$}


\newpage



\at(0cm,0cm){\includegraphics[width=8cm,bb=0 0 1920 1080]{./youtube/thumbnails/templates/smart_background/三角関数.jpeg}}
{\color{orange}\bf\boldmath\huge\underline{三角関数の定義}}\vspace{0.5zw}

\large
\bf\boldmath 問．下の各角$\theta$に対して，

\huge 
\vspace{-0.1zw}
\hspace{0.1zw}
$\cos \theta ,\;\sin \theta ,\;\tan \theta$
% \vspace{-1zw}
\vspace{0.2zw}


\large
\hfill の値を求めよ．
\vspace{0.7zw}

\small
(1)  $\bunsuu{5}{6}\pi$\hfill
(2)  $\bunsuu{5}{4}\pi$\hfill
(3)  $-\bunsuu{\pi }{3}$\hfill
(4)  $3\pi$\hfill
(5)  $-\bunsuu{{11}}{6}\pi$

\at(7.0cm,0.2cm){\small\color{bradorange}$\overset{\text{三角関数}}{\text{基本}}$}


\newpage



\at(0cm,0cm){\includegraphics[width=8cm,bb=0 0 1920 1080]{./youtube/thumbnails/templates/smart_background/三角関数.jpeg}}
{\color{orange}\bf\boldmath\huge\underline{弧度法とは？}}\vspace{0.3zw}

\Large 
\bf\boldmath 問．半径が$r,\;$中心角が$\theta$の扇形の弧の長さを$l,\;$面積を$S$とすると，

\huge
\vspace{0.3zw}
\hspace{0.5zw}$l=r\theta ,\;S=\bunsuu{1}{2}r^2\theta \vspace{0.3zw}$


\at(7.0cm,0.2cm){\small\color{bradorange}$\overset{\text{三角関数}}{\text{基本}}$}


\newpage



\at(0cm,0cm){\includegraphics[width=8cm,bb=0 0 1920 1080]{./youtube/thumbnails/templates/smart_background/三角関数.jpeg}}
{\color{orange}\bf\boldmath\huge\underline{一般角,動径とは？}}\vspace{0.3zw}

\large 
\bf\boldmath 問．動径$\text{OP}$と始線$\text{OX}$のなす角の$1$つを$\alpha$とすると，動径$\text{OP}$の表す一般角は，

\LARGE
\hspace{0.2zw}
$\alpha +{360}^\circ ×n \hspace{0.5zw}(n$は整数$)$

\large 
\vspace{0.5zw}
\hfill と表せる．
\at(7.0cm,0.2cm){\small\color{bradorange}$\overset{\text{三角関数}}{\text{基本}}$}


\end{document}

